%% ===========================================================================
%%  Manuscript for Climatic Change (Springer Nature)   sn-jnl, sn-basic
%%  Compile with:  pdflatex main / bibtex main / pdflatex main / pdflatex main
%%
%%  DOUBLE-BLIND: this file must contain no author names, affiliations,
%%  emails, ORCIDs, acknowledgements, funding text, or identifying URLs.
%%  All of that lives in title_page.tex, which is NOT input by this file.
%% ===========================================================================
\documentclass[pdflatex,sn-basic,lineno]{sn-jnl}

\usepackage{amsmath,amssymb}
\usepackage{graphicx}
\usepackage{booktabs}
\usepackage{array}
\usepackage{url}

\graphicspath{{figures/}}

\newcommand{\degC}{\ensuremath{^{\circ}\mathrm{C}}}
\newcommand{\ph}[1]{\textbf{[PLACEHOLDER: #1]}}
\newcommand{\vi}{\textbf{[verify-intended]}}

\begin{document}

%% -- Title --
\title{Projections of Heat-Stress Conditions Across Bangladesh Using Explainable Machine Learning and Climate Trend Integration}%Where and How Fast Is Heat Stress Worsening? A District-Level, Multi-Index Machine-Learning Projection for 64 districts of Bangladesh }

%% -- Authors and affiliations are deliberately absent: see title_page.tex --

%% -- Abstract --
\abstract{Rising heat and humidity increasingly threaten health, labour, and energy systems
across Bangladesh, yet operational planning needs district-resolved,
uncertainty-aware projections of the specific conditions that drive human heat
stress. We present a deterministic, reproducible machine-learning framework that
projects seven temperature-correlated heat-stress conditions, heat-index comfort,
wet-bulb temperature, indoor and outdoor wet-bulb globe temperature, cooling
degree-days, ultraviolet exposure, and dew-point moisture stress
for all 64
districts of Bangladesh. Four gradient-boosting ensembles are trained on the
2014--2024 monthly record using future-known seasonal features; the observed
1980--2024 Theil--Sen trend is superimposed so that projections evolve
realistically, and all thermal indices are recomputed from the projected primary
variables to avoid target leakage. Validated against withheld 2025 observations,
the projected indices attain median $R^{2}$ of $0.95$--$0.98$, and their
operational risk classes match observations to within one class in $100\%$ of
district-months. A Mann-Kendall analysis shows the maximum temperature rising in
every district (median $+0.56\,^{\circ}$C per decade) while minimum temperature
remains flat, indicating a widening diurnal range; the coastal and southwestern
zones carry the greatest combined burden, and ultraviolet and oppressive humidity
hazards are effectively nationwide. District choropleths, zonal envelopes, and
calibrated uncertainty bands provide a directly actionable heat-stress outlook
while transparently bounding the models' skill relative to climatology. The
complete 64-district, seven-condition projection is openly explorable through an
accompanying interactive web application.}

%% -- Keywords --
\keywords{Temperature-Correlated Conditions, public health, Bangladesh, seasonal projection, Applied Machine Learning}

\maketitle

\end{document}
